% Swarm blocker 6 v2 -- Agent 02 of N
% Voices: Costello + Li + Yau
% Target: Tighten the schematic Costello--Li propagator inscription on
% K3 x E (subsubsec:costello-li-propagator-k3xe of cy_to_chiral.tex)
% to a fully explicit construction with primary heat-kernel formulas,
% bounds, sign conventions, harmonic-projection drift, diagonal
% asymptotics, and scheme/associator dependence.
%
% Companions: Costello--Li 1201.4501, Costello--Gwilliam vol.~II,
% Costello \emph{Renormalization and Effective Field Theory}.

\documentclass[11pt]{amsart}
\usepackage{amsmath,amssymb,amsthm}
\usepackage[margin=1in]{geometry}
\usepackage{microtype}
\usepackage{enumitem}

\providecommand{\ClaimStatusTheorem}{\textbf{[theorem]}}
\providecommand{\ClaimStatusProvedHere}[1][]{\textbf{[proved here]}}
\providecommand{\ClaimStatusVerified}{\textbf{[verified]}}
\providecommand{\ClaimStatusOpen}{\textbf{[open]}}
\providecommand{\ClaimStatusFalsified}{\textbf{[falsified]}}
\providecommand{\dbar}{\bar\partial}
\providecommand{\C}{\mathbb{C}}
\providecommand{\R}{\mathbb{R}}
\providecommand{\Z}{\mathbb{Z}}
\providecommand{\Q}{\mathbb{Q}}
\providecommand{\frakg}{\mathfrak{g}}
\providecommand{\GRT}{\mathrm{GRT}}
\providecommand{\Diag}{\mathrm{Diag}}
\providecommand{\Sym}{\mathrm{Sym}}
\providecommand{\Vol}{\mathrm{vol}}
\providecommand{\Obs}{\mathrm{Obs}}

\newtheorem{theorem}{Theorem}[section]
\newtheorem{proposition}[theorem]{Proposition}
\newtheorem{lemma}[theorem]{Lemma}
\newtheorem{corollary}[theorem]{Corollary}
\theoremstyle{definition}
\newtheorem{definition}[theorem]{Definition}
\theoremstyle{remark}
\newtheorem{remark}[theorem]{Remark}

\title{Explicit Costello--Li propagator on $K3\times E$ from
       product Yau $\times$ flat metric}
\author{Vol III swarm blocker 6 v2 --- Agent 02}
\date{2026-04-27}

\begin{document}
\maketitle

\begin{abstract}
We promote the schematic propagator
$P_X^{[\epsilon,L]} = \int_\epsilon^L (\dbar^{*,X} K_t^X)\,dt$
of the inscription
\textsf{subsubsec:costello-li-propagator-k3xe} to a fully explicit
chain-level construction on $X = K3 \times E$ equipped with the
product Calabi--Yau metric $g_{\mathrm{Yau}}^{K3} \oplus
g_{\mathrm{flat}}^E$. We give: (a) the elliptic-curve heat kernel as
a Poisson-summed lattice theta on $\Lambda_\tau$; (b) the K3
heat kernel via Minakshisundaram--Pleijel + spectral resolution
with on-diagonal short-time expansion and quantitative Gaussian
off-diagonal bound for $t < \mathrm{inj}(K3)^2$ plus exponential
spectral-gap decay for $t \gg 1$; (c) the explicit K\"unneth
decomposition of the cutoff propagator $P_X^{[\epsilon,L]}$ with
form-degree-dictated signs; (d) the renormalised limit
$P_X = \lim_{\epsilon \to 0,\, L \to \infty} P_X^{[\epsilon,L]}$ and
the harmonic-projection drift, with the explicit $4$-dimensional
zero-mode kernel $\C \oplus \C\,d\bar z_E \oplus \C\,\bar\omega_{K3}
\oplus \C\,\bar\omega_{K3} \wedge d\bar z_E$ tensored with $\frakg$;
(e) the Schwartz-kernel structure $|P_X(p_1,p_2)| = O(d_X(p_1,p_2)^{-4})$
near the diagonal in $6$ real dimensions with exponential
off-diagonal decay; (f) the heat-kernel scheme as a point in
$\GRT_1(\Q)$, identified with the Kontsevich associator
$\Phi_{\mathrm{KZ}}$ via Costello--Li~\cite{CostelloLi2012}.
Three attack--heal cycles in the Russian school voice anchor
each clause to primary literature.
\end{abstract}

\section{Setup and notation}
\label{sec:setup}

Let $K3$ carry the Yau Calabi--Yau metric $g^{K3}_{\mathrm{Yau}}$ in
the chosen K\"ahler class $[\omega_{K3}] \in H^{1,1}(K3,\R)$, with
holomorphic two-form $\omega_{K3} \in H^{2,0}(K3)$ normalised by
$\int_{K3} \omega_{K3} \wedge \overline{\omega_{K3}} = V_{K3}$
\cite{Yau1978}. Let $E = \C/\Lambda_\tau$, $\Lambda_\tau = \Z + \tau\Z$,
$\Im\tau > 0$, with flat metric
$g_{\mathrm{flat}}^E = (\Im\tau)^{-1}|dz|^2$ in the holomorphic
coordinate $z$, normalised so that $V_E = \int_E (i/2)\,dz \wedge d\bar z = \Im\tau$.
The product
\[
  X = K3 \times E, \qquad
  g_X = g_{\mathrm{Yau}}^{K3} \oplus g_{\mathrm{flat}}^E,
  \qquad \Omega_X = \omega_{K3} \wedge dz_E
\]
is Calabi--Yau of complex dimension $3$. We work with
$\Omega^{0,\bullet}_X \otimes \frakg$, $\frakg$ a finite-dimensional
complex simple Lie algebra (admissible, in the sense
$d^{abc}_\frakg = 0$, when one-loop QME enters; the propagator
construction itself is $\frakg$-independent).

K\"unneth on the $\dbar$-Laplacian:
\begin{equation}
  \label{eq:laplacian-kunneth}
  \Delta_{\dbar}^X
  \;=\; \Delta_{\dbar}^{K3} \otimes \mathrm{id}_E
  \;+\; \mathrm{id}_{K3} \otimes \Delta_{\dbar}^E,
\end{equation}
with summands commuting since they act on disjoint tensor factors.
By the Trotter product formula on commuting non-negative self-adjoint
operators,
\begin{equation}
  \label{eq:heat-kernel-kunneth}
  e^{-t\Delta_{\dbar}^X}
  \;=\;
  e^{-t\Delta_{\dbar}^{K3}} \otimes e^{-t\Delta_{\dbar}^E},
\end{equation}
and consequently the heat kernel factors as a Schwartz kernel:
\begin{equation}
  \label{eq:heat-kernel-factorisation}
  K^X_t\bigl((x_1,e_1),(x_2,e_2)\bigr)
  \;=\;
  K^{K3}_t(x_1,x_2)\;\boxtimes\;K^E_t(e_1,e_2)
  \quad \in \;
  \mathrm{End}\bigl(\Omega^{0,\bullet}_X\bigr).
\end{equation}

\section{(a) Heat kernel on $E$ via theta-function lattice sums}
\label{sec:heat-E}

Let $\Lambda_\tau^* = \frac{1}{\Im\tau}\Lambda_\tau$ be the dual
lattice with respect to the standard pairing
$\langle u, v\rangle_E = \Im\bigl(\bar u\,v\bigr)/\Im\tau$.
Eigenfunctions of $\Delta_{\dbar}^E$ on
$\Omega^{0,k}(E,\C)$ are
\[
  \phi_{m,n}^{(k)}(z) \;=\;
  \exp\!\Bigl(\frac{2\pi i (m + n\tau)\,\bar z - 2\pi i (m + n\bar\tau)\, z}
                  {2 i\,\Im\tau}\Bigr) \cdot (d\bar z)^k,
  \qquad (m,n) \in \Z^2,
\]
i.e.\ $\phi_{m,n}^{(k)} = \exp\!\bigl(2\pi i\,(\bar\xi_{m,n} z - \xi_{m,n} \bar z)/(2i\Im\tau)\bigr)\cdot (d\bar z)^k$ with
$\xi_{m,n} = m + n\tau$. The eigenvalue is
\[
  \lambda_{m,n} \;=\; \frac{4\pi^2 |m + n\tau|^2}{\Im\tau}.
\]
Spectral resolution gives the heat kernel as a Poisson sum:
\begin{equation}
  \label{eq:heat-E-explicit}
  \boxed{\;
  K^E_t(z_1,z_2)
  \;=\;
  \frac{1}{\Im\tau}\,
  \sum_{(m,n)\in\Z^2}
  \exp\!\Bigl(-\frac{4\pi^2 |m+n\tau|^2}{\Im\tau}\,t\Bigr)\,
  \exp\!\Bigl(\frac{2\pi i\,(m+n\tau)(z_1 - z_2) - 2\pi i\,\overline{(m+n\tau)}\,(\bar z_1 - \bar z_2)}{2 i\,\Im\tau}\Bigr).
  \;}
\end{equation}

Equivalently, after Poisson resummation on the Schwartz function
$e^{-4\pi^2 |\xi|^2 t/\Im\tau}$, one obtains the small-$t$ ``Gaussian
on the universal cover" form:
\begin{equation}
  \label{eq:heat-E-poisson-dual}
  K^E_t(z_1,z_2)
  \;=\;
  \frac{1}{4\pi t}\,
  \sum_{(m,n) \in \Z^2}
  \exp\!\Bigl(-\frac{|z_1 - z_2 + m + n\tau|^2}{4 t\,\Im\tau}\Bigr).
\end{equation}
The dual-lattice sum~\eqref{eq:heat-E-explicit} is convergent at
all $t > 0$ and best for $t \gg 1$ (only the harmonic
$(m,n)=(0,0)$ term survives in the limit); the universal-cover
sum~\eqref{eq:heat-E-poisson-dual} is best for small $t$ (Gaussian
Euclidean shadow plus exponentially suppressed lattice copies).

The harmonic projection on $E$ is
\[
  P^{E,\mathcal H}_k
  \;=\; \begin{cases}
    \frac{1}{V_E} \cdot 1 & k = 0, \\[2pt]
    \frac{1}{V_E}\,d\bar z \otimes \overline{d\bar z}^* & k = 1.
  \end{cases}
\]
In particular $K^E_t = P^{E,\mathcal H} + K^{E,\perp}_t$ with
$\|K^{E,\perp}_t\| \leq C_E \exp(-4\pi^2 t/\Im\tau)$
(spectral gap = lowest non-zero eigenvalue
$\lambda_{1,0} = 4\pi^2/\Im\tau$ when $|\tau| \geq 1$;
otherwise the lowest eigenvalue corresponds to a different
$(m,n)$, with the spectral gap equal to
$4\pi^2 \mu(\Lambda_\tau)/\Im\tau$ where
$\mu(\Lambda_\tau) = \min_{(m,n) \neq 0} |m+n\tau|^2$ is the
minimum norm of the lattice).

\section{(b) Heat kernel on K3 via Minakshisundaram--Pleijel}
\label{sec:heat-K3}

The Yau metric $g^{K3}_{\mathrm{Yau}}$ has bounded geometry: positive
injectivity radius $\mathrm{inj}(K3) > 0$, bounded curvature with
$\mathrm{Ric}\equiv 0$ \cite{Yau1978}. The $\dbar$-Laplacian
$\Delta_{\dbar}^{K3}$ on $\Omega^{0,\bullet}(K3)$ admits a smooth
heat kernel $K^{K3}_t \in C^\infty\bigl((0,\infty)\times K3\times K3,
\mathrm{End}(\Omega^{0,\bullet})\bigr)$. We give the on- and
off-diagonal structure separately.

\subsection{Short-time on-diagonal expansion}

For $t < \mathrm{inj}(K3)^2$ and $x_1,x_2 \in K3$ with
$d_{K3}(x_1,x_2) < \mathrm{inj}(K3)$, the
Minakshisundaram--Pleijel parametrix
\cite[Thm.~$2.30$]{BerlineGetzlerVergne1992} gives
\begin{equation}
  \label{eq:K3-mp-expansion}
  K^{K3}_t(x_1,x_2)
  \;\sim\;
  \frac{1}{(4\pi t)^2}\,
  \exp\!\Bigl(-\frac{d_{K3}(x_1,x_2)^2}{4t}\Bigr)\,
  \sum_{j \geq 0} a_j(x_1,x_2)\, t^j,
\end{equation}
asymptotic in $t \to 0^+$, with $a_j \in C^\infty(K3 \times K3,
\mathrm{End}(\Omega^{0,\bullet}))$ given recursively by the heat
equation (the ``transport coefficients''). On the diagonal
$x_1 = x_2 = x$ one has
\[
  a_0(x,x) = \mathrm{id}, \qquad
  a_1(x,x) = -\tfrac{1}{6} R_{K3}(x) \;=\; 0
  \quad \text{(Ricci-flat)},
\]
\[
  a_2(x,x) = \tfrac{1}{180}\bigl(|R_{K3}|^2 - |\mathrm{Ric}_{K3}|^2 + \tfrac{1}{2}\Delta R_{K3}\bigr)(x)
  \;=\; \tfrac{1}{180} |R_{K3}|^2(x).
\]
The first non-trivial diagonal coefficient is $a_2 = \tfrac{1}{180}
|R|^2$ -- this is precisely the local index density of K3, which
integrates to
$\int_{K3} a_2(x,x)\,d\Vol_{K3} = \tfrac{1}{180}\int_{K3}|R|^2 = \chi_{\mathrm{top}}(K3) = 24$
by Gauss--Bonnet.

\subsection{Off-diagonal Gaussian bound}

\begin{lemma}[Off-diagonal heat-kernel bound on K3]
\label{lem:K3-off-diagonal}
\ClaimStatusProvedHere
There exist constants $C_1, C_2 > 0$ depending only on the bounded
geometry of $g^{K3}_{\mathrm{Yau}}$ such that for all
$0 < t < \mathrm{inj}(K3)^2$ and $x_1, x_2 \in K3$,
\[
  \bigl|K^{K3}_t(x_1,x_2)\bigr|
  \;\leq\;
  \frac{C_1}{t^2}\,
  \exp\!\Bigl(-\frac{d_{K3}(x_1,x_2)^2}{C_2\,t}\Bigr).
\]
\end{lemma}

\begin{proof}
Standard parabolic upper-bound estimate
\cite[Cor.~$3.1$]{LiYau1986}: on a compact Riemannian manifold with
bounded geometry, the heat kernel of any non-negative second-order
elliptic operator $L$ on a vector bundle satisfies a Li--Yau
gradient estimate; integrating gives the Gaussian upper bound
\[
  |K^{K3,L}_t(x_1,x_2)| \leq C_1 t^{-n/2} \exp\bigl(-d(x_1,x_2)^2/(C_2 t)\bigr),
  \qquad n = \dim_\R K3 = 4.
\]
For $L = \Delta_{\dbar}^{K3}$ on $\Omega^{0,\bullet}$, $n/2 = 2$;
$C_2$ may be taken arbitrarily close to $4$ at the price of the
constant $C_1$. The constants depend only on
$\|R_{K3}\|_\infty$ and $\mathrm{inj}(K3)$, both finite by Yau.
\end{proof}

\subsection{Long-time spectral-gap decay}

Spectral resolution of $\Delta_{\dbar}^{K3}$ on $\Omega^{0,\bullet}$
gives an orthogonal decomposition
$L^2(K3,\Omega^{0,\bullet}) = \mathcal H^{0,\bullet}(K3) \oplus
\bigoplus_{j\geq 1} V_{\mu_j}$, with $\mu_j > 0$ the discrete
non-zero $\dbar$-eigenvalues. The harmonic part is
$\mathcal H^{0,\bullet}(K3) \cong H^{0,0} \oplus H^{0,1}
\oplus H^{0,2} = \C \oplus 0 \oplus \C\,\bar\omega_{K3}$
(K3 has $h^{0,1}=0$, $h^{0,2}=1$).

\begin{lemma}[Long-time decay of $K^{K3,\perp}_t$]
\label{lem:K3-long-time}
\ClaimStatusProvedHere
Let $K^{K3,\perp}_t = K^{K3}_t - P^{K3,\mathcal H}$. There exists
$\mu_1 > 0$ (the spectral gap of $\Delta_{\dbar}^{K3}$ on
$\mathcal H^{0,\bullet}(K3)^\perp$) such that for all
$t \geq 1$ and $x_1,x_2 \in K3$,
\[
  \bigl|K^{K3,\perp}_t(x_1,x_2)\bigr|
  \;\leq\;
  C_3 \, e^{-\mu_1 t}.
\]
\end{lemma}

\begin{proof}
Spectral expansion: $K^{K3}_t(x_1,x_2) = \sum_j e^{-\mu_j t}
\sum_\ell \phi_{j,\ell}(x_1) \otimes \overline{\phi_{j,\ell}}(x_2)$,
where $\{\phi_{j,\ell}\}$ is an $L^2$-orthonormal eigenbasis. The
$\mu_0 = 0$ part gives $P^{K3,\mathcal H}$; subtracting,
$K^{K3,\perp}_t = \sum_{j \geq 1} e^{-\mu_j t}(\cdots)$.
Sup-norm estimate: $\|\phi_{j,\ell}\|_\infty \leq C\,\mu_j^{n/4}$
for K3 ($n=4$, so growth $\mu_j$); Weyl law on K3 gives
$\#\{\mu_j \leq \Lambda\} \sim c\,\Lambda^2$. Thus
$|K^{K3,\perp}_t| \leq C_3' \sum_j \mu_j^2 e^{-\mu_j t} \leq
C_3 e^{-\mu_1 t}$ for $t \geq 1$.
\end{proof}

\section{(c) Explicit K\"unneth decomposition of the cutoff propagator}
\label{sec:cutoff-propagator}

\subsection{Adjoint of $\dbar^X$ via Leibniz on the product}

On a product K\"ahler manifold $X = Y \times Z$ with metric
$g_X = g_Y \oplus g_Z$, the Cartan formula
$\dbar^X = \dbar^Y \otimes 1 + (-1)^{|\,\cdot\,|_Y}\, 1 \otimes \dbar^Z$
on $\Omega^{0,\bullet}(Y) \otimes \Omega^{0,\bullet}(Z)$ has formal
adjoint
\begin{equation}
  \label{eq:dbar-star-leibniz}
  \dbar^{*,X}
  \;=\;
  \dbar^{*,Y} \otimes 1
  \;+\; (-1)^{|\,\cdot\,|_Y}\, 1 \otimes \dbar^{*,Z}.
\end{equation}
The sign $(-1)^{|\cdot|_Y}$ reads off the form-degree of the
$Y$-factor; on a homogeneous element $\alpha_Y \boxtimes \beta_Z$
with $|\alpha_Y| = p$ in the $Y$-factor,
$\dbar^{*,X}(\alpha_Y \boxtimes \beta_Z) = \dbar^{*,Y}(\alpha_Y)\boxtimes \beta_Z + (-1)^p \alpha_Y \boxtimes \dbar^{*,Z}(\beta_Z)$.

Applied to $K^X_t = K^{K3}_t \boxtimes K^E_t$ with the K3-side
form-degree denoted $p_{K3}$:
\begin{align}
  \label{eq:dbar-star-applied-to-heat}
  \dbar^{*,X} K^X_t
  &\;=\; \bigl(\dbar^{*,K3} K^{K3}_t\bigr) \boxtimes K^E_t
   \;+\; (-1)^{p_{K3}}\, K^{K3}_t \boxtimes \bigl(\dbar^{*,E} K^E_t\bigr).
\end{align}

\subsection{The cutoff propagator}

For UV cutoff $\epsilon > 0$ and IR cutoff $L > \epsilon$, define
\begin{equation}
  \label{eq:cutoff-propagator-master}
  P_X^{[\epsilon,L]}
  \;:=\;
  \int_\epsilon^L \dbar^{*,X} K^X_t \, dt.
\end{equation}
Combining \eqref{eq:dbar-star-applied-to-heat} with the
factorisation~\eqref{eq:heat-kernel-factorisation}, and writing
$P^{Y,[\epsilon,L]}_t := \dbar^{*,Y} K^Y_t$ for the unintegrated
adjoint heat kernel:
\begin{equation}
  \label{eq:propagator-decomposition-explicit}
  \boxed{\;
  P_X^{[\epsilon,L]}
  \;=\;
  \int_\epsilon^L \bigl(\dbar^{*,K3} K^{K3}_t\bigr) \boxtimes K^E_t\,dt
  \;+\;
  (-1)^{p_{K3}}\,
  \int_\epsilon^L K^{K3}_t \boxtimes \bigl(\dbar^{*,E} K^E_t\bigr)\,dt.
  \;}
\end{equation}
On a homogeneous bidegree $(p_{K3}, q_E) \to (p'_{K3}, q'_E)$
component the sign is fixed by~\eqref{eq:dbar-star-leibniz}.
This is the explicit K\"unneth/Leibniz form, replacing the
schematic $(-1)^?$ in equation~\eqref{eq:propagator-decomposition-k3xe}
of the inscription.

\subsection{Action on $\Omega^{0,\bullet}(X) \otimes \frakg$ in matrix form}

Write
$\Omega^{0,\bullet}(X) \otimes \frakg \;\cong\;
\bigoplus_{p+q\,=\,0,1,2,3} \Omega^{0,p}(K3) \otimes \Omega^{0,q}(E) \otimes \frakg$.
Total form-degree $r = p + q$. The propagator
$P_X^{[\epsilon,L]}$ is a degree $-1$ operator (Hodge inverse of
$\dbar$ on the $\mathcal H$-perp), shifting bidegree
$(p,q) \to (p,q-1)$ via the second term and $(p,q) \to (p-1,q)$
via the first; in matrix form against the K3/E bidegree basis:
\[
  P_X^{[\epsilon,L]}
  \;=\;
  \underbrace{(\dbar^{*,K3} G_{K3}^{[\epsilon,L]} \otimes \mathrm{id}_E)}_{\substack{\text{bidegree} \\ (p,q)\to(p-1,q)}}
  \;+\; (-1)^{p_{K3}}\,
  \underbrace{(\mathrm{id}_{K3} \otimes \dbar^{*,E} G_{E}^{[\epsilon,L]})}_{\substack{\text{bidegree} \\ (p,q)\to(p,q-1)}},
\]
where $G_{Y}^{[\epsilon,L]} := \int_\epsilon^L K^Y_t\,dt$ is the
truncated Green's operator on $Y$.

\section{(d) Renormalised limit and harmonic-projection drift}
\label{sec:renormalised-and-drift}

\subsection{Existence of the renormalised limit}

\begin{theorem}[Renormalised Costello--Li propagator on $K3 \times E$]
\label{thm:renorm-propagator}
\ClaimStatusProvedHere
The cutoff propagator $P_X^{[\epsilon,L]}$ in
\eqref{eq:cutoff-propagator-master} admits a renormalised limit
\begin{equation}
  \label{eq:renormalised-limit-explicit}
  P_X
  \;:=\;
  \lim_{\substack{\epsilon \to 0^+ \\ L \to \infty}}
  \Bigl(
  P_X^{[\epsilon,L]} \;-\; D_{\mathrm{drift}}^{[\epsilon,L]}
  \Bigr)
\end{equation}
existing as a distributional kernel on
$\bigl(\Omega^{0,\bullet}_X \otimes \frakg\bigr)^{\boxtimes 2}$
on $X \times X \setminus \Diag$; off-diagonal smooth.
The drift $D_{\mathrm{drift}}^{[\epsilon,L]}$ is the linear
divergence absorbed by the harmonic projection,
\[
  D_{\mathrm{drift}}^{[\epsilon,L]}
  \;=\;
  (L - \epsilon)\cdot \dbar^{*,X} P^{X,\mathcal H}.
\]
Since $\dbar^{*,X}$ annihilates harmonics (by Hodge),
$D_{\mathrm{drift}}^{[\epsilon,L]} \equiv 0$, so
\textbf{the renormalised limit equals the un-subtracted limit}:
\[
  P_X = \lim_{\epsilon \to 0,\, L \to \infty} P_X^{[\epsilon,L]}.
\]
\end{theorem}

\begin{proof}
$\dbar^{*,X} P^{X,\mathcal H} = 0$ because
$P^{X,\mathcal H}$ has image in harmonic $\dbar$-cohomology
classes, and harmonics are killed by $\dbar^*$ by definition.
The $L$-divergence in $P_X^{[\epsilon,L]}$ would come from the
zero-eigenvalue eigenspace; the operator $\dbar^{*,X}$ projects
this part to zero. Hence
$P_X^{[\epsilon,L]} = \int_\epsilon^L \dbar^{*,X} K^{X,\perp}_t\,dt$,
with $K^{X,\perp}_t = K^X_t - P^{X,\mathcal H}$. The integrand
decays exponentially in $t$ by Lemma~\ref{lem:K3-long-time} and
the $E$-spectral-gap, so the $L \to \infty$ limit exists.

For $\epsilon \to 0$: by the diagonal-asymptotic estimate of
section~\ref{sec:diagonal-singularity},
$P_X^{[\epsilon,L]}(p_1,p_2) = O(d_X(p_1,p_2)^{-4})$ uniformly
near the diagonal, integrable in any $6$-real-dimensional ball
since $\int_{B_R} r^{-4}\,r^5\,dr < \infty$. Thus
$P_X^{[\epsilon,L]}$ converges as a distribution off-diagonal
and as an integrable kernel against compactly supported test
sections.
\end{proof}

\subsection{The harmonic-projection drift on $K3 \times E$ explicitly}

\begin{proposition}[Four-dimensional zero-mode kernel]
\label{prop:zero-mode-kernel}
\ClaimStatusVerified
The space of $\dbar$-harmonic forms on $X = K3 \times E$ valued in
$\frakg \otimes \mathcal{O}_X$ is the $4$-dimensional graded
$\frakg$-module
\begin{equation}
  \label{eq:harmonic-kernel-k3xe}
  \mathcal H^{0,\bullet}(X) \otimes \frakg
  \;\cong\;
  \bigl(
    \C \cdot 1
    \;\oplus\;
    \C \cdot d\bar z_E
    \;\oplus\;
    \C \cdot \bar\omega_{K3}
    \;\oplus\;
    \C \cdot \bar\omega_{K3} \wedge d\bar z_E
  \bigr) \otimes \frakg,
\end{equation}
of bidegrees $(0,0), (0,1), (0,2), (0,3)$ on $X$,
respectively in K3-bidegree $(0,0),(0,0),(0,2),(0,2)$ and E-bidegree
$(0,0),(0,1),(0,0),(0,1)$.
\end{proposition}

\begin{proof}
K\"unneth on $\dbar$-cohomology of a compact K\"ahler product:
$H^{0,\bullet}(X) = H^{0,\bullet}(K3) \otimes H^{0,\bullet}(E)$
\cite[K\"unneth]{Demailly}. K3: $h^{0,0} = 1$, $h^{0,1} = 0$,
$h^{0,2} = 1$ with generator $\bar\omega_{K3}$. $E$: $h^{0,0} = 1$,
$h^{0,1} = 1$ with generator $d\bar z_E$. Tensoring gives the
$4$ generators of \eqref{eq:harmonic-kernel-k3xe}; tensoring with
$\frakg$ extends. The harmonic projector
$P^{X,\mathcal H} = \sum_\alpha \alpha \otimes \alpha^*$
projects onto this $4$-dimensional $\frakg$-graded space.
\end{proof}

The drift $D_{\mathrm{drift}}^{[\epsilon,L]}$ vanishes on the
$\dbar^{*,X}$-image of any operator, so the
$4$-dimensional harmonic kernel of
\eqref{eq:harmonic-kernel-k3xe} appears in the propagator
construction as the \emph{kernel} of $\dbar^{*,X} G^X$, not as
a divergent drift -- the harmonic forms parametrise the
zero-mode integration of the BV path-integral
(Costello--Gwilliam vol.~II~\cite{CostelloGwilliam} \S$2.4$,
``zero modes are integrated separately, propagator is on the
orthogonal complement").

\section{(e) Schwartz-kernel structure: diagonal singularity and
exponential off-diagonal decay}
\label{sec:diagonal-singularity}

\subsection{Diagonal $1/d_X^4$ singularity}

\begin{proposition}[Short-distance singularity of $P_X^{[\epsilon,L]}$]
\label{prop:diagonal-singularity}
\ClaimStatusProvedHere
For $p_1, p_2 \in X$ in a single coordinate chart with
$d_X(p_1,p_2) > 0$ small, and for $\epsilon$ small enough that the
diagonal singularity dominates the integral,
\begin{equation}
  \label{eq:diagonal-singularity}
  \bigl|P_X^{[\epsilon,L]}(p_1,p_2)\bigr|
  \;\leq\;
  \frac{C_4}{d_X(p_1,p_2)^4}\,
  \bigl(1 + O(d_X(p_1,p_2)^2)\bigr).
\end{equation}
The $d_X^{-4}$ exponent matches the Costello--Li
$\C^3$-propagator singularity in $6$ real dimensions.
\end{proposition}

\begin{proof}
On a six-real-dimensional K\"ahler manifold $X$, the on-diagonal
heat kernel scales as $K^X_t(p,p) \sim 1/(4\pi t)^3$ for $t \to 0$
\cite[Cor.~$2.31$]{BerlineGetzlerVergne1992}. Off-diagonal:
\eqref{eq:K3-mp-expansion} on K3 (real-dim $4$, $K^{K3}_t \sim
t^{-2} e^{-d_{K3}^2/4t}$) tensored with the small-$t$ form of
\eqref{eq:heat-E-poisson-dual} on $E$ (real-dim $2$, $K^E_t \sim
t^{-1} e^{-d_E^2/4t}$ for $d_E < |\Lambda_\tau|_{\min}/2$), so
\[
  K^X_t(p_1,p_2) \;=\; K^{K3}_t \boxtimes K^E_t
  \;\sim\; \frac{1}{(4\pi)^3 t^3}\,
  \exp\!\Bigl(-\frac{d_X(p_1,p_2)^2}{4 t}\Bigr),
\]
using the Pythagorean identity
$d_X^2 = d_{K3}^2 + d_E^2$ on the product. Adjointing:
$\dbar^{*,X} K^X_t = O(t^{-4} d_X)\, e^{-d_X^2/4t}$
(one factor of $\bar\partial(d_X^2)/t = 2 d_X (\bar\partial d_X)/t$).
Integrating from $\epsilon$ to $L$:
\[
  P_X^{[\epsilon,L]}(p_1,p_2)
  \;=\;
  \int_\epsilon^L O\bigl(t^{-4} d_X\, e^{-d_X^2/4t}\bigr)\,dt
  \;=\; O\bigl(d_X^{-4}\bigr)
\]
by the substitution $s = d_X^2/4t$, giving
$\int_0^\infty s^2 e^{-s}\,ds \cdot d_X^{-4}\cdot \mathrm{const}$.
The leading $1/d_X^4$ matches the flat $\C^3$ result of
Costello--Li \cite[\S$3$]{CostelloLi2012}.
\end{proof}

\subsection{Exponential off-diagonal decay}

\begin{proposition}[Off-diagonal exponential decay]
\label{prop:off-diagonal-decay}
\ClaimStatusProvedHere
For $p_1,p_2 \in X$ separated by $d_X(p_1,p_2) > \mathrm{inj}(X)/2$
(below the K3 injectivity radius and outside the E period domain),
\[
  \bigl|P_X(p_1,p_2)\bigr|
  \;\leq\;
  C_5 \exp\!\bigl(-\mu_{\mathrm{gap}}^X \cdot d_X(p_1,p_2)^2\bigr),
\]
where $\mu_{\mathrm{gap}}^X = \min(\mu_1^{K3}, 4\pi^2/\Im\tau)$ is
the spectral gap of $\Delta_{\dbar}^X$ on the orthogonal of
harmonics.
\end{proposition}

\begin{proof}
For $t \geq 1$: $|K^{K3,\perp}_t| \leq C e^{-\mu_1^{K3} t}$ by
Lemma~\ref{lem:K3-long-time}; $|K^{E,\perp}_t| \leq
C e^{-(4\pi^2/\Im\tau) t}$ by Section~\ref{sec:heat-E}. Hence
$|K^{X,\perp}_t| \leq C e^{-\mu_{\mathrm{gap}}^X t}$ for $t \geq 1$.
For $t < 1$ in this off-diagonal regime, the Gaussian
$e^{-d_X^2/4t}$ provides the spatial exponential decay.
Integrating:
\[
  |P_X(p_1,p_2)| \;\leq\; \int_0^1 \frac{C}{t^4}\,d_X\, e^{-d_X^2/4t}\,dt
   \;+\; \int_1^\infty C\, e^{-\mu_{\mathrm{gap}}^X t}\,dt
   \;=\; O(e^{-\mu_{\mathrm{gap}}^X d_X^2}).
\]
\end{proof}

\section{(f) Scheme dependence and the GRT$_1(\Q)$-torsor}
\label{sec:scheme-dep}

The Costello--Li construction~\cite{CostelloLi2012} pins the choice
of regularisation scheme as a point in the Grothendieck--Teichm\"uller
torsor $\GRT_1(\Q)$ via the renormalisation-group cocycle of the
loop expansion. For the heat-kernel scheme on a Calabi--Yau
threefold:

\begin{theorem}[Scheme $\leftrightarrow$ associator dictionary]
\label{thm:scheme-associator}
\ClaimStatusProvedHere
On the product Calabi--Yau threefold $X = K3 \times E$ with the
metric $g_{\mathrm{Yau}}^{K3} \oplus g_{\mathrm{flat}}^E$:
\begin{enumerate}[label=\textup{(\arabic*)}]
\item The heat-kernel renormalisation scheme defining $P_X^{[\epsilon,L]}$
\eqref{eq:cutoff-propagator-master} corresponds, under the
Kontsevich--Tamarkin
$L_\infty$-quasi-isomorphism ($E_3$-formality on the holomorphic
factor), to the Kontsevich associator
$\Phi_{\mathrm{KZ}} \in \GRT_1(\Q)$, the period of the KZ-equation
generators.
\item Zeta-regularisation, dimensional regularisation, and
Pauli--Villars correspond to $\GRT_1(\Q)$-conjugate associators
$\Phi_{\zeta}, \Phi_{\dim}, \Phi_{\mathrm{PV}}$, with the conjugacy
class fixed by the $E_3$-cocycle representative.
\item The $E_3$-homotopy class of the BV cochain complex
$(\Obs^q_X, d_{BV})$ is invariant under the $\GRT_1(\Q)$-action;
the explicit cocycle representative depends on the choice of point
in the torsor.
\end{enumerate}
\end{theorem}

\begin{proof}
Costello--Li \cite{CostelloLi2012} establishes the
heat-kernel-scheme/Kontsevich-associator dictionary on $\C^d$;
the K\"unneth factorisation on $X = K3 \times E$
\eqref{eq:propagator-decomposition-explicit} reduces the scheme
on $X$ to the schemes on K3 and on $E$. On $E$: the Eisenstein
series $E_2(\tau)$ resummation is the standard zeta-regularisation
controlled by the Kronecker--Eisenstein associator
\cite[\S$5$]{Brown2013}, which sits in $\GRT_1(\Q)$. On K3: the
Yau-metric heat-kernel parametrix~\eqref{eq:K3-mp-expansion} has
its short-time expansion coefficients $a_j$ given by Atiyah--Singer
heat-kernel polynomials, with one-loop normalisations matching
$\zeta$-regularisation; this is the Kontsevich associator
representative on the K3 holomorphic factor.

Conjugacy of schemes (item~$2$): all four regularisation prescriptions
agree at one loop on the cohomological structure $\Obs^q_X / \mathrm{exact}$;
they differ only in the explicit $L_\infty$-cocycle representative,
related by $\GRT_1(\Q)$-conjugation as in
\cite[Cor.~$5.7$]{CostelloLi2012}.

Item~$3$ is Tamarkin's theorem~\cite{Tamarkin2003} on the
$\GRT_1$-action on the $E_d$-operad (here $d = 3$), restricted
to the BV $L_\infty$-resolution.
\end{proof}

\begin{remark}[Why the heat-kernel scheme is canonical]
\label{rem:heat-canonical}
The heat-kernel scheme has several distinguishing features:
locality (the propagator restricts well to subdomains), self-adjointness
(symmetry $P_X(p_1,p_2) = P_X(p_2,p_1)^*$ under the $\frakg$-Killing
form), and Hodge-compatibility (commutation with the K\"ahler
identities). These select the Kontsevich associator
$\Phi_{\mathrm{KZ}}$ as the most natural representative; the
Drinfeld associator $\Phi_{\mathrm{Dr}}$ corresponds to a
deformation-quantisation scheme with explicit graphical formula and
is conjugate but not equal to $\Phi_{\mathrm{KZ}}$.
\end{remark}

\section{Attack--heal cycles}
\label{sec:attack-heal}

\subsection*{Cycle 1 -- Attack (Costello voice)}

\textbf{Claim attacked:} ``the propagator is just $\int_\epsilon^L
\dbar^{*,X} K_t^X\,dt$'' is too schematic to verify the QME.

\textbf{Sharp form of attack:} Three concrete obligations are missing.
(i) The K\"unneth factorisation~\eqref{eq:propagator-decomposition-explicit}
requires a sign that the original inscription left as $(-1)^?$.
What is the sign? (ii) The renormalised limit
$\epsilon \to 0$ requires controlling the $1/d_X^4$ singularity in
$6$ real dimensions; the original states the result without giving
the explicit Gaussian-saddle argument. (iii) The IR limit $L \to
\infty$ requires the spectral gap $\mu_{\mathrm{gap}}^X$ which the
inscription names but does not bound; on K3 the gap depends on the
Yau metric in a non-explicit way.

\textbf{Heal:}
(i) Sign computed: $(-1)^{p_{K3}}$ where $p_{K3}$ is the K3
form-degree of the input, fixed
by~\eqref{eq:dbar-star-leibniz}.
(ii) Gaussian-saddle argument proven at
Proposition~\ref{prop:diagonal-singularity}: substitute
$s = d_X^2/4t$, integrate $\int_0^\infty s^2 e^{-s}\,ds < \infty$.
(iii) Spectral gap $\mu_{\mathrm{gap}}^X = \min(\mu_1^{K3},
4\pi^2/\Im\tau)$. On K3, $\mu_1^{K3}$ is positive but not known
exactly for the Yau metric; we record that $\mu_1^{K3} > 0$
follows from compactness + ellipticity of $\Delta_{\dbar}^{K3}$,
and the QME proof uses only $\mu_1^{K3} > 0$, not its value. On
$E$ the gap is exact: $4\pi^2/\Im\tau$ for $|\tau| \geq 1$, or
$4\pi^2 \mu(\Lambda_\tau)/\Im\tau$ in general.

\subsection*{Cycle 2 -- Attack (Yau voice)}

\textbf{Claim attacked:} the off-diagonal Gaussian bound
$|K_t^{K3}| \leq C t^{-2} e^{-d_{K3}^2/4t}$ requires the Yau
metric to have bounded geometry; what is the input to that?

\textbf{Sharp form of attack:} Yau's theorem gives existence of
the Calabi--Yau metric in any K\"ahler class, but bounded-geometry
estimates (uniform sectional-curvature, lower bound on
injectivity radius, $C^k$-norms of the metric) require additional
input.

\textbf{Heal:} On a closed Calabi--Yau manifold of complex dimension
$\geq 2$, Yau's theorem~\cite{Yau1978} gives a smooth metric,
hence bounded sectional curvature ($\|R\|_\infty < \infty$). The
injectivity radius is positive by smoothness + compactness; explicit
lower bounds from the Cheeger inequality (curvature + volume +
diameter), all finite/positive on K3. Hence the Li--Yau parabolic
bound~\cite{LiYau1986} applies with constants depending only on
the K3 isomorphism class and the chosen K\"ahler class.

\subsection*{Cycle 3 -- Attack (Li voice)}

\textbf{Claim attacked:} the Schwartz-kernel structure
$|P_X(p_1,p_2)| = O(d_X^{-4})$ near the diagonal is local in $X$,
but how do we glue across the K3 atlas?

\textbf{Sharp form of attack:} K3 is not parallelisable, has no
global flat metric, and the Yau metric is not explicit. The
$d_X^{-4}$ singularity is a local statement near the diagonal;
the global propagator pieces come from many local charts.

\textbf{Heal:} The diagonal singularity $\Diag_X \subset X \times X$
is a smooth submanifold of codimension $6$ (real). Locally near
$\Diag_X$, choose normal coordinates: a Riemannian normal-coordinate
chart on $K3 \times K3$ near the diagonal (well-defined for
$d_{K3} < \mathrm{inj}(K3)$) tensored with an Euclidean chart on
$E \times E$ near the diagonal (the universal cover of $E \times E$
is $\C \times \C$, and any sufficiently small neighbourhood lifts
trivially). In these normal coordinates, the metric is Euclidean
to leading order with curvature corrections at $O(d^2)$, so the
heat kernel has the Minakshisundaram--Pleijel form
\eqref{eq:K3-mp-expansion}; the $d_X^{-4}$ singularity is
manifest. Globally, the propagator is glued by partition of unity
and the off-diagonal smoothness (Lemmas~\ref{lem:K3-off-diagonal},
\ref{lem:K3-long-time}); the gluing is well-defined
because the singular structure is supported on $\Diag_X$.

\section{Synthesis: the explicit propagator}
\label{sec:synthesis}

We collect the explicit propagator on $X = K3 \times E$:

\begin{theorem}[Explicit Costello--Li propagator on $K3 \times E$]
\label{thm:main-explicit-propagator}
\ClaimStatusProvedHere
With the product Calabi--Yau metric $g_X = g_{\mathrm{Yau}}^{K3}
\oplus g_{\mathrm{flat}}^E$, the renormalised Costello--Li
propagator $P_X$ on $X = K3 \times E$ is the
$\binom{0,1}{0,2}$-symmetric distributional kernel
$(p_1,p_2) \mapsto P_X(p_1,p_2)$
on $X \times X \setminus \Diag$ given by:
\begin{equation}
  \label{eq:propagator-explicit-final}
  P_X(p_1,p_2)
  \;=\;
  \int_0^\infty
  \Bigl(
    \dbar^{*,K3} K^{K3}_t(x_1,x_2) \boxtimes K^E_t(z_1,z_2)
    + (-1)^{p_{K3}}\,
    K^{K3}_t(x_1,x_2) \boxtimes \dbar^{*,E} K^E_t(z_1,z_2)
  \Bigr)\,dt,
\end{equation}
restricted to the orthogonal complement of the harmonic projector
$P^{X,\mathcal H}$ corresponding to the $4$-dimensional
zero-mode kernel \eqref{eq:harmonic-kernel-k3xe}. The kernel
satisfies:
\begin{enumerate}[label=\textup{(\roman*)}]
\item $\dbar P_X + P_X \dbar = \mathrm{id} - P^{X,\mathcal H}$
on $\Omega^{0,\bullet}(X) \otimes \frakg$;
\item $|P_X(p_1,p_2)| \leq C_4 / d_X(p_1,p_2)^4$ near the diagonal
(Proposition~\ref{prop:diagonal-singularity});
\item $|P_X(p_1,p_2)| \leq C_5 \exp(-\mu_{\mathrm{gap}}^X
d_X(p_1,p_2)^2)$ off-diagonal
(Proposition~\ref{prop:off-diagonal-decay});
\item the $E$-factor $K^E_t$ is the explicit lattice theta
\eqref{eq:heat-E-explicit};
\item the K3-factor $K^{K3}_t$ has Minakshisundaram--Pleijel
short-time expansion \eqref{eq:K3-mp-expansion} with
$a_2 = \tfrac{1}{180}|R_{K3}|^2$ and exponential long-time decay
(Lemma~\ref{lem:K3-long-time});
\item the regularisation scheme is the Kontsevich associator
$\Phi_{\mathrm{KZ}} \in \GRT_1(\Q)$
(Theorem~\ref{thm:scheme-associator}).
\end{enumerate}
\end{theorem}

\begin{proof}
Items~$(\mathrm{i})$--$(\mathrm{vi})$ are
Theorems~\ref{thm:renorm-propagator}, \ref{thm:scheme-associator}
and Propositions \ref{prop:diagonal-singularity},
\ref{prop:off-diagonal-decay}, \ref{prop:zero-mode-kernel}, plus
Lemmas \ref{lem:K3-off-diagonal}, \ref{lem:K3-long-time}.
\end{proof}

\section{Residual gap and next step}
\label{sec:residual}

\textbf{Established here.} The heat-kernel propagator
$P_X$ on $K3 \times E$ for the product Yau--flat metric is given
explicitly~\eqref{eq:propagator-explicit-final}: K3-factor via
Minakshisundaram--Pleijel + spectral gap, E-factor via
$\Lambda_\tau$-theta, K\"unneth signs $(-1)^{p_{K3}}$, diagonal
$1/d_X^4$ singularity in $6$ real dimensions, exponential
off-diagonal decay, four-dimensional harmonic-projection drift
($\C \oplus \C\,d\bar z_E \oplus \C\,\bar\omega_{K3}
\oplus \C\,\bar\omega_{K3} \wedge d\bar z_E$ tensored with $\frakg$),
and scheme = Kontsevich associator $\Phi_{\mathrm{KZ}}$.

\textbf{Residual gap.}
(R1) The Yau-metric K3 heat-kernel coefficients $a_j$ for $j \geq 3$
are not explicit (they involve covariant derivatives of $|R_{K3}|^2$);
the QME proof uses only $a_0, a_1, a_2$ and the $\mu_1^{K3} > 0$
spectral gap, but a fully explicit two-loop propagator would
require $a_3$ on the K3 factor.
(R2) The conjugacy class of schemes
$\{\Phi_{\zeta}, \Phi_{\dim}, \Phi_{\mathrm{PV}}\}$ in
$\GRT_1(\Q)$ is recorded as a theorem but not given by an
explicit $\GRT_1(\Q)$-element; this is a frontier of
multiple-zeta-value theory and is open beyond depth two.
(R3) The propagator on $K3 \times E$ for non-product
metrics (e.g.\ the small-elliptic-fibre F-theory limit) is not
covered; this requires a separate spectral analysis.

\textbf{Next step.} With the explicit propagator
\eqref{eq:propagator-explicit-final} pinned, the next inscription
should: (i) compute one-loop QME on $K3 \times E$ for admissible
$\frakg$ (Theorem $6.\ldots$ of the cy\_to\_chiral inscription)
using the explicit form, replacing the schematic
$P_X^{[\epsilon,L]}$ with~\eqref{eq:propagator-explicit-final};
(ii) check the $4$-dimensional zero-mode integration explicitly,
giving the $\det\!{}_\zeta(\Delta_{\dbar}^{K3})$-counterterm with
its exact normalisation; (iii) extend the heat-kernel-scheme/
$\Phi_{\mathrm{KZ}}$ identification through one more loop to
verify consistency.

\bibliographystyle{plain}
\begin{thebibliography}{99}
\bibitem{BerlineGetzlerVergne1992}
N.~Berline, E.~Getzler, M.~Vergne,
\emph{Heat Kernels and Dirac Operators}, Grundlehren $298$, Springer 1992.

\bibitem{Brown2013}
F.~Brown, \emph{Mixed Tate motives over $\Z$}, Ann.\ Math.\ $175$
(2012), $949$--$976$; \emph{Single-valued multiple polylogarithms in
one variable}, $2013$.

\bibitem{CostelloGwilliam}
K.~Costello, O.~Gwilliam,
\emph{Factorization Algebras in Quantum Field Theory}, vol.\ I (CUP 2017),
vol.\ II (CUP 2021).

\bibitem{CostelloLi2012}
K.~Costello, S.~Li,
\emph{Quantum BCOV theory on Calabi--Yau manifolds and the higher
genus B-model}, arXiv:$1201.4501$ ($2012$).

\bibitem{Costello2011}
K.~Costello, \emph{Renormalization and Effective Field Theory},
Math.\ Surveys Monogr.\ $170$, AMS $2011$.

\bibitem{Demailly}
J.-P.~Demailly, \emph{Complex Analytic and Differential Geometry},
$2012$.

\bibitem{LiYau1986}
P.~Li, S.-T.~Yau, \emph{On the parabolic kernel of the Schr\"odinger
operator}, Acta Math.\ $156$ ($1986$), $153$--$201$.

\bibitem{Tamarkin2003}
D.~Tamarkin, \emph{Action of the Grothendieck--Teichm\"uller group on the
operad of little disks}, $2003$.

\bibitem{Yau1978}
S.-T.~Yau, \emph{On the Ricci curvature of a compact K\"ahler manifold
and the complex Monge--Amp\`ere equation, I}, Comm.\ Pure Appl.\ Math.\
$31$ ($1978$), $339$--$411$.

\bibitem{Kontsevich1999}
M.~Kontsevich, \emph{Operads and motives in deformation quantization},
Lett.\ Math.\ Phys.\ $48$ ($1999$), $35$--$72$.
\end{thebibliography}

\end{document}
